\section{INTRODUCTION}

In 2113 the expanding Terran Confederation reaches the edge of the Ziru
Sirka, the First Imperium. Neither polity can accept the other's idea of
order. What begins as a frontier dispute becomes a succession of short,
violent wars separated by uneasy peaces. Terra learns quickly; the Imperial
provincial governor commands greater resources but must answer to a distant
Emperor.

\textit{Space Empires: Imperium} is a two-player strategic campaign. It keeps
the point-to-point jump map, political asymmetry, Glory system, and repeated
wars of GDW's \Imperium. It uses the ship groups, firing classes, attack and
defense values, hull damage, screening, fighters, and ground forces from
\SpaceEmpires\ and \CloseEncounters. Its Ship Yards make the primary Worlds
the vulnerable industrial centers of each side. A shared card deck introduces
a small amount of operational surprise.

Fleet combat specifically reconciles the two parent systems rather than
discarding either one. Fleets contest Long or Close Range and choose beam,
missile, or mixed armament in the manner of \Imperium, while attacks use the
printed classes, d10 Hit Numbers, screening, and hull damage of
\SpaceEmpires. Traveller-inspired pickets, point defense, and sandcasters
make the range and weapon choices visible at fleet scale.

This is an alpha playtest design. Its systems are complete enough to play the
\textit{First Shots} war, but the force mix and economy have not yet been
balanced through repeated play.

\subsection{The Central Design Choice}

The game is not a full 4X game. The map is known, the frontier is already
settled, and there are no non-player aliens. Players fight over an existing
political geography. Technology changes slowly and only between wars. The
interesting decisions should be where to build a logistics chain, whether to
hold missile range or close to beam range, when to commit the fleet, whether
to risk an invasion before orbital combat is settled, and how much political
capital to spend for a short-term advantage.

\subsection{Rules Hierarchy}

Use these rules first. If they do not address a situation, use the following
sources in order:

\begin{enumerate}
  \item this rulebook and the text in its card manifest;
  \item \SpaceEmpires\ rules for ships, groups, damage, screening, fighters,
    and retreat;
  \item \CloseEncounters\ section 14 for transports and ground combat; and
  \item \Imperium\ only for interpretation of its map and named systems.
\end{enumerate}

Rules from a parent game that concern hex movement, exploration, colony
growth, aliens, or that game's sequence of play do not apply. When printed
card text disagrees with the manifest in this rulebook, the manifest controls.

\subsection{Rounding and Dice}

Round fractions down unless a rule says otherwise. A ``die'' or ``d10'' is a
ten-sided die; a zero face counts as 10. ``2d6'' means two six-sided dice
added together.
